\begin{statement}[C.6.T2; Proved, with exact number-field local arithmetic]
Put $s=\sqrt3$ and
$t_\pm=(3008\pm1728s)/11$. Over $K=\Q(s)$, a global minimal
model for the fibre at $t_-$ is
\[
 E_-:y^2=x^3-(18+10s)x+21+12s,\qquad P_-=(0,3+2s).
\]
The model and point for $t_+$ are its conjugates. The
discriminant of this model is $1032912+596352s$. Define
\[
 \mathfrak p_2=(2,1+s),\qquad\mathfrak p_3=(s),\qquad
 \mathfrak p_{11,-}=(11,s+5),\quad
 \mathfrak p_{11,+}=(11,s-5).
\]
The bad-prime data of $E_-$ are
\[
\begin{array}{c|c|c|c|c|c}
v&e_v&v(\Delta)&\text{type}&\text{component order of }P_-&\lambda_v(P_-)\\\hline
\mathfrak p_2&2&8&\mathrm{IV}^*&3&0\\
\mathfrak p_3&2&3&\mathrm{III}&2&0\\
\mathfrak p_{11,-}&1&1&\mathrm I_1&1&\log11/12.
\end{array}
\]
For $E_+$ replace $\mathfrak p_{11,-}$ by
$\mathfrak p_{11,+}$. The nonzero component at
$\mathfrak p_3$ is unique; the two nonzero classes at
$\mathfrak p_2$ are interchanged by $P\mapsto-P$ and the
class of $P$ is a generator. All omitted local heights vanish.
Both points have order six, and
\begin{align*}
 \log N(P_-)&=\frac1{12}\log(17+10\sqrt3),\\
 \log N(P_+)&=\frac1{12}\log(10\sqrt3-17),\\
 \sum_{\sigma:K\hookrightarrow\R}\log N(P_\sigma)
 &=\frac{\log11}{12}
 =\sum_{v<\infty}n_v\lambda_v(P)-2\widehat h(P).
\end{align*}
The common value in the last line is
\[
 0.199824606066530878671828631497094108318475571161451431268214\ldots.
\]
In particular $A_\mp=(17+10\sqrt3)^{1/12}$ and
$(10\sqrt3-17)^{1/12}$ respectively give
$D_{\rm part}(t_\mp)=(6/\pi)\log(A_\mp/(2\sqrt3))$.

The natural Mordell--Weil group here is $E_\pm(K)$; the
notation $E_{t_\pm}(\Q)$ requested in the prompt does not
define an elliptic curve over $\Q$, since its $j$-invariant
is irrational. Over $K$, both curves have rank one and torsion
group $\Z/6\Z$. Thus the index of the torsion section in the
free quotient is infinite. For example $(1,-1-s)$ has infinite
order on $E_-$.
\end{statement}
\begin{proof}
For $t_-$ the change from the original family to this model
has scale $u=48(-64+37s)/11$, with translation parameters zero.
It is obtained by multiplying the first global minimalization
scale by the unit $(2+s)^{-3}$. Hence it preserves all finite
minimal valuations. The conjugate scale gives $E_+$.
Tate's algorithm yields the displayed table. Its discriminant
ideal is
$\mathfrak p_2^8\mathfrak p_3^3\mathfrak p_{11,-}$ and its
conductor ideal is
$\mathfrak p_2^2\mathfrak p_3^2\mathfrak p_{11,-}$.

As $5P=-P$, the normalized local heights are computed exactly by
\[
 \frac{\lambda_v(P)}{\log p}
 =\frac1{e_v}\left(\frac{v(\Delta)}{12}
                         -\frac{v(\psi_5(P))}{24}\right).
\]
The valuations of $\psi_5(P)$ at the three listed primes are
$16,6,0$, respectively. This proves the table, including zero
contribution at the two additive primes. Here $f_v=1$ at all
three primes, so weighting by $n_v=e_vf_v$ gives precisely
$\log11/12$. The exact torsion values proved in C.4 have product
$(17+10\sqrt3)(10\sqrt3-17)=11$, proving the archimedean side.

For the rank, translate the rational two-torsion point to zero.
The resulting equation is
\[
 y^2=x^3-(9+6s)x^2+(45+26s)x.
\]
The $2$-isogeny descent gives
$[E(K):\widehat\phi E'(K)]=2$ and
$[E'(K):\phi E(K)]=4$. Every surviving covering has a
displayed point; every excluded covering has a recorded local
obstruction at $\mathfrak p_2$ or $\mathfrak p_3$.
The quotient formula $2^{r+2}=2\cdot4$ proves $r=1$.
The field arithmetic was certified by \texttt{bnfcertify};
Simon descent was run with \texttt{LIMBIGPRIME=0}, so all local
tests were deterministic. The complete transcript is
\path{code/Crank_torsion_descent_sage.out}. The unmodified
Simon scripts maintained in Sage are retained with their source
URLs; the older upstream scripts and their unsuccessful
compatibility attempts are also retained. The torsion calculation
gives $\Z/6\Z$, generated by $P$, independently of the rank
calculation.
\end{proof}
